% Lösungen Einheit 5 von 5 – Prozentuale Veränderung (zusammenbau v0.1)
\einheitenkopf{Einheit 5 von 5 · Prozentuale Veränderung}

\erg{23}{a) $2{,}40\,€$}
%% TODO Lösung der Erklärzeile 24g) fehlt in der Bank
\erg{24}{a) auf $80\,\%$ gesenkt \quad b) $10\,\% = 8\,€$; $80 + 8 = 88\,€$ \quad c) $25\,\% = 15$ kg; $60 - 15 = 45$ kg \quad d) $5\,\% = 20$; $400 + 20 = 420$ Einwohner \quad e) $20\,\% = 10\,€$; $50 - 10 = 40\,€$ \quad f) $50\,\% = 45$ cm; $90 + 45 = 135$ cm \quad g) \ldots \quad h) $45 \cdot 1{,}2 = 54\,€$ \quad i) $10 : 40 = 0{,}25 = 25\,\%$ \quad j) $360 : 0{,}8 = 450\,€$ \quad k) $200 \cdot 1{,}19 = 238\,€$ \quad l) $6$ Prozentpunkte; $6 : 58 \approx 0{,}103 \approx 10{,}3\,\%$ mehr \quad m) $14 : 200 = 0{,}07 = 7\,\%$ \quad n) $1{,}29 - 1{,}15 = 0{,}14$; $0{,}14 : 1{,}15 \approx 0{,}122 \approx 12{,}2\,\%$ \quad o) $2{,}40 \cdot 1{,}25 = 3{,}00\,€$ (nicht nur $0{,}60\,€$) \quad p) voller Preis $2 \cdot 14 + 2 \cdot 9 = 46\,€$; $46 \cdot 0{,}85 = 39{,}10\,€$ (nicht der Rabatt $6{,}90\,€$) \quad q) Lücken: $100$ m und $5$ m; $15 : 240 = 0{,}0625 \approx 6{,}3\,\% > 5\,\%$ – Herr Lenz hat recht.}
\erg{25}{a) Fehler: Ida teilt durch den neuen Preis. Richtig: $15 : 60 = 0{,}25 = 25\,\%$. \quad b) Die Veränderung wird vom Ausgangswert aus gemessen; der alte Preis ist das Ganze, die $100\,\%$. \quad c) A: $600 \cdot 0{,}9 \cdot 0{,}9 = 486\,€$; B: $600 \cdot 0{,}8 = 480\,€$ – Angebot B ist $486 - 480 = 6\,€$ günstiger.}
